\chapteropen{03}

\begin{cele}
  \item prowadzić audyt czasu przez siedem dni na kartce albo w prostym arkuszu, z kategoriami i minutami, tak aby suma tygodnia dała dokładnie 10\,080 minut;
  \item rozpoznać pięć najczęstszych pożeraczy czasu w dobie dorosłego człowieka i odróżnić je od czasu, który nie podlega cięciu (sen, praca, ruch, jedzenie);
  \item policzyć, ile godzin odzyskujesz zmianą i ile kosztuje samo przejście, a z tych dwóch liczb wyprowadzić czas zwrotu;
  \item ocenić, czy zasada Pareto działa w twoim tygodniu, czy czas jest rozłożony równomiernie na wiele drobnych pozycji;
  \item rozpoznać pułapkę nieskończonej optymalizacji i wiedzieć, kiedy audyt przestaje być narzędziem, a zaczyna być formą odkładania.
\end{cele}

\section{Dlaczego audyt, a nie plan}

Większość ludzi, którzy chcą odzyskać czas, zaczyna od planu. Siadają w niedzielę wieczorem, rysują siatkę godzin od szóstej do dwudziestej drugiej i wpisują w nią to, co powinno się wydarzyć: trening, gotowanie, czytanie, pięć minut na medytację. Plan jest przyjemny w pisaniu, bo opisuje dzień, którego jeszcze nie ma. Problem polega na tym, że wszystko, czego w nim brakuje, zostaje w cieniu.

Brakuje w nim głównie tego, czego nie uważamy za wybór. Przeglądanie telefonu w łóżku nie wygląda na decyzję, tylko na chwilę przerwy. Dojazd nie wygląda na zajęcie, tylko na konieczność. Zakupy są po prostu czymś, co trzeba zrobić. Kiedy zapytać się wprost, ile czasu zjadło każde z tych zajęć, odpowiedzi są zwykle inne niż te, które podsuwa pamięć. \verify{Ludzie zwykle niedoszacowują czas spędzany w telefonie i przy ekranach}. Nie jest to zarzut wobec pamięci, tylko opis tego, jak działa uwaga: rzeczy nawykowe nie zostawiają śladu.

Dlatego ten rozdział nie zaczyna się od pytania, co wpisać w plan, tylko od pytania, dokąd naprawdę poszedł tydzień. Plan można napisać dopiero wtedy, gdy wiesz, z czego go budujesz.

\section{Jak prowadzić audyt przez siedem dni}

Audyt trwa siedem dni, a nie jeden, z trzech powodów. Po pierwsze, tydzień ma dwa rytmy: dni robocze i weekend, które różnią się tak mocno, że jeden dzień nie reprezentuje żadnego z nich. Po drugie, w pierwszym dniu prowadzenie zapisu zmienia zachowanie, więc pierwsze wpisy są najmniej wiarygodne. Po trzecie, siedem dni to jedyny okres, który jest równy tygodniowemu szablonowi, z którego później zbudujesz plan. Wszystkie liczby w tym rozdziale są tygodniowe z tego samego powodu.

Potrzebujesz dwóch narzędzi, i wystarczy jednego z nich.

\begin{itemize}
  \item \textbf{Kartka.} Wydrukuj tabelę z wierszami co 30 minut, od godziny wstania do godziny zaśnięcia, i wpisuj przy każdym wierszu kategorię i to, co faktycznie robiłeś. Kartka ma tę zaletę, że jest widoczna i wymusza zapis w chwili zmiany zajęcia, a nie wieczorem z pamięci. Jej wada: liczenie minut zajmuje czas i trudno policzyć sumy bez kalkulatora.
  \item \textbf{Arkusz.} Jeden wiersz to jedno zajęcie: data, godzina początku, godzina końca, kategoria, minuty (wzór $=$ koniec $-$ początek, po przeliczeniu na minuty), uwagi. Sumy po kategorii i po dniu liczy funkcja sumy warunkowej, a nie ręka. Jego wada: łatwo przestać go prowadzić w trzecim dniu.
\end{itemize}

Zapisuj każde zajęcie, które trwało co najmniej dziesięć minut. Krótsze zajęcia zbiorcze, na przykład sprawdzenie telefonu w windzie, zapisuj jako jeden wiersz z minutami w sumie, nie jako osobne wpisy. Dokładność na poziomie kwadransa wystarcza do decyzji i jest realistyczna do utrzymania; dokładność co do minuty kosztuje więcej uwagi niż sam audyt jest wart.

Kategorie ustal przed pierwszym dniem, nie w trakcie. Dziesięć, może dwanaście kategorii to wystarczająco dużo, żeby zobaczyć wzór, i na tyle mało, żeby zapis nie wymagał decyzji przy każdym wpisie. Sen, praca, dojazd, posiłki, ruch, sprzątanie i sprawy domowe, zakupy, zebrania i sprawy organizacyjne, ekrany rozrywkowe i reszta czasu wolnego to rozsądny punkt wyjścia. Jeżeli nie wiesz, do której kategorii włożyć dane zajęcie, wpisz je w osobną kategorię ,,nie wiem'' i sprawdź w niedzielę. Im więcej ,,nie wiem'' w arkuszu, tym mniej wiarygodny jest wynik.

Pilnuj jednej zasady: suma siedmiu dni musi wynosić $7 \cdot 24 \cdot 60 = 10\,080$ minut. Jeśli suma jest mniejsza, brakuje wpisów, a braki zwykle są w tym samym miejscu: w porach, których nie chcesz oglądać. Brakujące minuty nie są neutralne. Najczęściej to właśnie one okazują się największym pożeraczem.

\section{Kategorie pożeraczy czasu}

W większości tygodni pięć rodzajów zajęć zjada więcej godzin, niż ludzie podają, a ich wspólną cechą nie jest to, że są złe. Są ciągłe, niewidoczne i nie mają wyraźnego końca. Oto one, z krótkim opisem tego, jak wyglądają w arkuszu.

\textbf{Ekrany.} Wszystkie zajęcia przy ekranie łączymy w jedną grupę na etapie audytu, a rozdzielamy dopiero w analizie. Trzeba je rozdzielić, bo praca na komputerze, seriale i przeglądanie wiadomości mają różne koszty. Seriale i gry są odpoczynkiem, który wybierasz świadomie. Przeglądanie bez celu jest czasem, którego nie pamiętasz.

\textbf{Dojazd.} Czas w drodze do pracy i z niej to koszt stały, który rzadko zostaje zakwestionowany, bo wydaje się nieunikniony. Często nie jest. Różnica między dojazdem jednym środkiem a drugim, między podróżą w godzinach szczytu a poza nimi, albo między dojazdem a pracą zdalną bywa wielokrotnością tego, co zyskujesz na jakiejkolwiek pojedynczej zmianie w domu. Dojazd najlepiej rozdzielić na kategorię ,,dojazd'' i nie mieszać go z zakupami czy zebraniami.

\textbf{Zakupy.} Zakupy mają swój koszt w sklepie, w kolejce i w drodze, a do tego dochodzi koszt decyzji: co kupić, czy to jest potrzebne, czego zabrakło. Jeżeli sklep jest po drodze, zakupy wydają się darmowe. Jeżeli wymagają oddzielnego wyjścia, każde wyjście jest pełnym kosztem.

\textbf{Scrollowanie.} To część ekranów, którą zasługuje na osobną pozycję, bo ma najkrótsze i najbardziej rozproszone zajęcia. Jedno sprawdzenie telefonu trwa minutę, pięć takich minut dziennie nie wygląda na coś, o czym warto mówić. Suma siedmiu dni potrafi jednak wynieść kilka godzin, a w arkuszu widać ją dopiero po zsumowaniu.

\textbf{Zebrania.} Zebrania w pracy są zwykle w kalendarzu i liczą się jako praca. Zebrania poza pracą, czyli rada rodziców, wspólnota mieszkaniowa, koło zainteresowań, sprawy urzędowe, spotkania rodzinne z planowaniem, rzadko są w kalendarzu jako zajęcie, i rzadko ktoś je liczy. Zapisuj je oddzielnie, bo ich suma bywa zaskakująca, a ich koszt rośnie wraz z liczbą osób, które w nich uczestniczą.

Obok pożeraczy są zajęcia, których nie wolno wycinać, i o nich w sekcji o cięciu. Rozróżnienie jest ważne: audyt, który traktuje sen i ruch jak pozycje do ograniczenia, prowadzi do planu, który nie da się utrzymać.

\section{Przykład: tydzień w tabeli}

Poniżej tydzień przykładowy. Liczby nie pochodzą z żadnego badania; to rachunek ilustrujący, jak sumuje się dobrze prowadzony audyt. Przyjęto dwanaście godzin snu każdej doby, ponieważ taki budżet jest punktem wyjścia tej książki. Sen w audycie się nie zmienia: nadal go zapisujesz, tylko go nie optymalizujesz.

\begin{przyklad}[Tydzień przykładowy]
Osoba pracuje na pełny etat w biurze, z dojazdem. Pierwsza kolumna to tygodniowa suma godzin, druga to średnia dzienna, liczona jako suma tygodnia podzielona przez siedem.

\begin{center}\small
\begin{tabular}{@{}lrr@{}}
\toprule
\textbf{Kategoria} & \textbf{Tydzień [h]} & \textbf{Dziennie [h]}\\
\midrule
Sen & 84,0 & 12,0\\
Praca płatna (5 dni po 8 h) & 40,0 & 5,7\\
Dojazd (5 dni po 1 h 12 min) & 6,0 & 0,9\\
Scrollowanie w telefonie & 7,5 & 1,1\\
Seriale i gry & 5,0 & 0,7\\
Zakupy (w sklepie i w kolejce) & 3,0 & 0,4\\
Gotowanie i posiłki & 8,0 & 1,1\\
Sprzątanie i sprawy domowe & 4,0 & 0,6\\
Zebrania poza pracą & 2,5 & 0,4\\
Ruch i higiena & 3,0 & 0,4\\
Chleb (czas aktywny) & 1,0 & 0,1\\
Wolne i nieplanowane & 4,0 & 0,6\\
\midrule
\textbf{Razem} & \textbf{168,0} & \textbf{24,0}\\
\bottomrule
\end{tabular}
\end{center}

Sprawdzenie sumy: $84{,}0 + 40{,}0 + 6{,}0 + 7{,}5 + 5{,}0 + 3{,}0 + 8{,}0 + 4{,}0 + 2{,}5 + 3{,}0 + 1{,}0 + 4{,}0 = 168{,}0$. To jest $7 \cdot 24$, czyli cały tydzień. Gdyby suma wyszła na przykład 165, brakowałoby trzech godzin i trzeba by szukać ich w zapisach, zanim cokolwiek się wyciągnie z tabeli.

Z tej tabeli wynikają dwie rzeczy. Po pierwsze, po odjęciu snu i pracy zostaje $168 - 84 - 40 = 44$ godziny na wszystko inne. Po drugie, w tych czterdziestu czterech godzinach jest $7{,}5 + 5{,}0 = 12{,}5$ godziny ekranów, czyli więcej niż cały czas przeznaczony na zakupy, zebrania i ruch razem wzięte.
\end{przyklad}

\subsection*{Czy działa zasada Pareto w czasie?}

Zasada Pareto mówi, że niewielka część przyczyn daje większość efektów. Stosowana do czasu oznacza: kilka kategorii powinno zjadać większość godzin, więc to je warto ciąć. Sprawdźmy to na tygodniu przykładowym, licząc tylko czterdzieści cztery godziny poza snem i pracą.

Największe dwie kategorie to gotowanie (8,0) i scrollowanie (7,5). Razem $15{,}5$ godziny, czyli $15{,}5 / 44 \approx 35\%$ tego, co zostaje. Dwie z dziesięciu kategorii zjadają jedną trzecią, a reszta jest rozrzucona. Gdy dołożysz trzecią największą, dojazd (6,0), dostajesz $21{,}5 / 44 \approx 49\%$, czyli niemal połowę.

To nie jest dramatyczna koncentracja, a w tygodniach większości ludzi bywa jeszcze mniejsza. Czas nie rozkłada się zwykle jak na obrazku z podręcznika: nie ma jednego pożeracza ważącego osiemdziesiąt procent. Jest wiele średnich pozycji, z których każda wydaje się niewinna. Wniosek praktyczny jest taki, że trzeba ciąć wybrane pozycje, a nie szukać jednej, która wszystko rozwiąże. Gdybyś po audycie wskazał tylko jedną kategorię, najpewniej wybrałbyś tę, która najbardziej irytuje, a nie tę, która zjada najwięcej.

\begin{falle}[Pułapka: audyt, który mierzy samopoczucie]
Najczęstszy błąd na etapie analizy to ocena każdej kategorii według tego, jak się z nią czujesz, zamiast według tego, ile zjada i co dostajesz w zamian. Dwie godziny czytania sprawiają dobre wrażenie, a dwie godziny scrollowania złe, ale z punktu widzenia tygodnia obie zajmują dokładnie dwie godziny. Ocena ma sens, gdy rozdzielasz odpoczynek, który wybrałeś, od czasu, który po prostu przepłynął. Ocena według samopoczucia prowadzi do cięcia tego, co lubisz, zamiast tego, co nie daje ci nic.
\end{falle}

\section{Rachunek: godzina odzyskana kontra koszt}

Odzyskana godzina nie jest tym samym co wolna godzina. Jeśli wycinasz scrollowanie, a zyskany czas wracasz do scrollowania w innej formie, nie zyskałeś nic. Dlatego rachunek trzeba zrobić na dwóch kolumnach: ile godzin tygodniowo zyskujesz i ile godzin kosztuje samo przejście na nowy układ.

Dla każdej zmiany potrzebne są trzy liczby. Pierwsza to oszczędność tygodniowa $S$ w godzinach, czyli różnica między czasem przed i po. Druga to koszt jednorazowy $C$: godziny, które wydasz na ustawienie zmiany, naukę nowego nawyku, przygotowanie list, przepisów, wtyczek, rozmowę z domownikami. Trzecia to koszt bieżący $O$, czyli czas, który zmiana zjada każdego tygodnia, a którego nie wliczyłeś w $S$. Czas zwrotu wynosi wtedy
\[
  n = \frac{C}{S - O}
\]
tygodni, a po $t$ tygodniach zysk netto to $t \cdot (S - O) - C$.

Wzór jest prosty, a jego najczęstsze nadużycie polega na pominięciu $C$. Zmiana, która oszczędza trzy godziny tygodniowo i kosztuje dziesięć godzin wdrożenia, zwraca się po ponad trzech tygodniach. Zmiana, która oszczędza trzy godziny i kosztuje pół godziny, zwraca się w mniej niż dzień. Obie wyglądają na podobne, dopóki nie policzysz $C$.

\begin{przyklad}[Trzy zmiany z tygodnia przykładowego]
Każda z trzech zmian celuje w inną pozycję z tabeli. Liczby $C$ są szacunkami autora tego przykładu, a nie pomiarami; w twoim tygodniu policz je sam, zapisując czas faktycznie poświęcony na ustawienie.

\begin{center}\small
\begin{tabular}{@{}lrrr@{}}
\toprule
\textbf{Zmiana} & \textbf{$S$ [h/tydz.]} & \textbf{$C$ [h]} & \textbf{Zwrot [tydz.]}\\
\midrule
Wyłączenie powiadomień i usunięcie aplikacji z ekranu głównego & 3,0 & 0,5 & 0,2\\
Zakupy online z listą stałą, dwa wyjścia zamiast czterech & 2,0 & 3,0 & 1,5\\
Gotowanie hurtowe: trzy przepisy w rotacji, jedna sesja w tygodniu & 3,0 & 4,0 & 1,3\\
\bottomrule
\end{tabular}
\end{center}

Zwrot liczony jest jak wyżej: $0{,}5/3{,}0 \approx 0{,}2$, $3{,}0/2{,}0 = 1{,}5$, $4{,}0/3{,}0 \approx 1{,}3$. Pierwsza zmiana zwraca się w około półtora dnia, co jest rzadkim przypadkiem wartym zrobienia od razu.

Gdy zrobisz wszystkie trzy, oszczędność wynosi $3{,}0 + 2{,}0 + 3{,}0 = 8{,}0$ godziny tygodniowo. Jednorazowy koszt to $0{,}5 + 3{,}0 + 4{,}0 = 7{,}5$ godziny. Dołóż dwie godziny na przejście: w dwóch pierwszych tygodniach gotowanie idzie wolniej, a zakupy online się mylą, co razem daje $C = 9{,}5$ godziny. Zwrot całości wynosi $9{,}5 / 8{,}0 \approx 1{,}2$ tygodnia. Po dwunastu tygodniach, gdyby zmiany się utrzymały, zysk netto to $12 \cdot 8{,}0 - 9{,}5 = 86{,}5$ godziny, czyli ponad dwa pełne dni roboczych.
\end{przyklad}

Ten wynik pokazuje też, czego nie należy pomijać: zmiany przez pierwsze tygodnie są droższe, niż wynika z oszczędności. To jest \emph{cena przejścia}, o której jest osobna sekcja. Najbardziej uczciwy sposób to policzyć ją z góry i sprawdzić po dwóch tygodniach, czy szacunek był bliski.

\begin{uwaga}[Zyskane godziny trzeba zarejestrować]
Odzyskany czas idzie gdzieś. Jeżeli oszczędności z tabeli wrócą do wolnego czasu, tabela po zmianach wygląda tak: scrollowanie 4,5, zakupy 1,0, gotowanie 5,0, a wolne i nieplanowane 12,0 zamiast 4,0. Suma nadal wynosi 168. Jeśli jednak po dwóch tygodniach wolne wróci do czterech godzin, a w arkuszu pojawią się nowe ekrany, zmiana nie odzyskała niczego, tylko przesunęła czas w inne miejsce. Zapisywanie po zmianie jest więc częścią zmiany, nie dodatkiem do niej.
\end{uwaga}

\section{Czego nie wycinać}

Audyt ma tendencję do wskazywania kandydatów na cięcie tam, gdzie nie powinno ich być. Istnieją zajęcia, które w rachunku wyglądają jak koszt, a w rzeczywistości są inwestycją, bez której reszta tygodnia się rozsypuje. Trzy z nich są nie do ruszenia w tym budżecie.

\textbf{Sen.} Dwanaście godzin dziennie to budżet, który ta książka zakłada, a nie cel, który audyt ma zredukować. Jeżeli w tabeli sen zjada mniej niż zakładasz, to nie jest pozycja do ścięcia, tylko sygnał, że trzeba zbadać, skąd się wzięła różnica. Sen w arkuszu zapisujesz tak samo jak wszystko inne, z godziną zaśnięcia i wstania, ale nie wolno go mierzyć jako przeszkody. Ile naprawdę potrzebujesz, sprawdza rozdział o śnie.

\textbf{Jedzenie.} Gotowanie i posiłki są w tabeli jednym z największych bloków, a mimo to nie należy ich wycinać, tylko organizować. Różnica między gotowaniem wieczorem co dzień a dwiema sesjami w tygodniu jest czasem, który można zdjąć z dnia, nie z talerza. Jedzenie samo w sobie jest częścią budżetu, którego nie chcesz ograniczać.

\textbf{Ruch.} Trzy godziny ruchu i higieny tygodniowo wyglądają w tabeli jak pozycja do przycięcia, ale ruch jest tym, co najłatwiej wypada z planu, gdy czas się kurczy, a jego brak wpływa na resztę dnia. \verify{Regularny ruch poprawia jakość snu u dorosłych}. Do tej pozycji stosuj zasadę odwrotną niż do ekranów: nie wycinaj, tylko dbaj, żeby miała stały czas w tygodniu.

Odpoczynek też nie jest kandydatem. Seriale i gry w tabeli to pięć godzin, które wybierasz. Ich cięcie bez zamiany prowadzi do zmęczenia, a zmęczenie wraca w postaci tej samej przerwy, tylko gorszej jakości. Ciąć można sposób, w jaki odpoczywasz, na przykład przechodząc z przeglądania na jeden wybrany serial, ale nie samą potrzebę odpoczynku.

\subsection*{Zamienniki zamiast pustki}

Cięcie bez zamiany zostawia lukę, którą wypełnia to, co było pod ręką, zwykle telefon. Dlatego każde cięcie w rachunku powinno mieć zamiennik wpisany z góry. Zamiennik nie musi być produktywny. Musi być zaplanowany, bo tylko zaplanowane zajęcie zachowuje się przy zmęczeniu. Jeśli wycinasz przeglądanie przed snem, zamiennikiem może być książka leżąca na stoliku nocnym. Jeśli wycinasz zakupy w sklepie, zamiennikiem jest zakupowa lista zapisana w telefonie i zamówienie raz w tygodniu.

\begin{zasada}[Wycinaj to, co ma najniższy koszt zamiany]
Nie wybieraj kategorii do cięcia według liczby godzin, tylko według stosunku zysku do kosztu przejścia. Dwie godziny w zebraniach, które dają się zlikwidować jednym mailem, są lepszym celem niż dwie godziny ekranów, które wymagają miesiąca przestawiania nawyków. Pierwsza zmiana zwraca się w dniach, druga w miesiącach.
\end{zasada}

\section{Cena przejścia}

Cena przejścia to koszt, który płacisz za samą zmianę, zanim zobaczysz jej efekt. Ma kilka składników i większość z nich jest niewidoczna w tabeli, bo nie jest czasem liczonym w minutach. Najczęściej są to: godziny ustawienia, czyli czas na konfigurację aplikacji, list, harmonogramów; okres gorszej sprawności, w którym nowe zajęcie idzie wolniej niż stare; opór domowników, którzy przyzwyczaili się do poprzedniego układu i trzeba ich przekonać albo przeczekać; oraz koszt pieniędzy, na przykład za dostawę zakupów, który nie jest w tym rozdziale liczony, ale w rachunku warto go zanotować osobno.

Najważniejsza jest pierwsza i trzecia pozycja, bo są wyraźnie większe, niż zakłada większość planów. Zmiana, która wygląda na dwie minuty dziennie, kosztuje w praktyce tydzień nauki, jeśli wymaga przestawienia wieczornej rutyny. Zmiana, która wymaga rozmowy z partnerem o podziale obowiązków, kosztuje tę rozmowę, a jej czas zależy od tego, jak wygląda ten podział dzisiaj.

Dlatego cenę przejścia trzeba wpisać do rachunku z góry, jako szacunek, a po dwóch tygodniach sprawdzić, czy był bliski. Jeśli był zaniżony dwukrotnie, zmiana nadal może się opłacać, tylko wolniej. Jeśli był zaniżony dziesięciokrotnie, to sygnał, że zmiana jest za duża i trzeba ją podzielić na mniejsze kroki, a nie upierać się przy całości.

\section{Pułapka: optymalizuję bez końca}

Audyt ma jedną cechę, którą łatwo pomylić z zaletą: daje poczucie kontroli. Każdy kolejny tydzień zapisu dostarcza nowych liczb, nowych kategorii, nowych wykresów. Można tak prowadzić arkusz miesiącami, dopracowując kategorie, dodając kolumny z nastrojem, liczbą kroków i godziną najlepszej koncentracji, a potem analizując to wszystko w kolejnym arkuszu. Wtedy audyt przestaje być narzędziem do podjęcia decyzji i staje się sposobem na odłożenie decyzji.

Miarą tej pułapki jest stosunek czasu poświęconego na pomiar do czasu odzyskanego przez zmiany. Jeśli audyt zabrał czternaście godzin, a zmiany odzyskały trzy godziny tygodniowo, zwrot następuje po pięciu tygodniach, ale tylko wtedy, gdy zmiany faktycznie wdrożysz. Jeśli arkusz rośnie, a zmian nie ma, koszt audytu rośnie, a zysk wynosi zero.

Ogranicz więc audyt z góry. Siedem dni, jeden arkusz, jedna analiza i lista co najwyżej trzech zmian. Po wdrożeniu tych trzech zmian nie dokładaj kolejnych kategorii, tylko zobacz, co przyniosły. Drugi audyt robisz po dwóch miesiącach, nie po dwóch tygodniach, i tylko wtedy, gdy wiesz, czego szukasz.

\begin{cwiczenia}
  \exitem{Ile minut dziennie to 7,5 godziny tygodniowo?}{$7{,}5 / 7 \approx 1{,}07$ godziny, czyli około $64$ minut dziennie. Liczba minut dziennych dla tygodniowej sumy to zawsze tygodniowa suma podzielona przez siedem, zanim zaczniesz z nią cokolwiek porównywać.}
  \exitem{Zsumuj tabelę tygodnia przykładowego i pokaż, że daje 168 godzin.}{$84{,}0 + 40{,}0 + 6{,}0 + 7{,}5 + 5{,}0 + 3{,}0 + 8{,}0 + 4{,}0 + 2{,}5 + 3{,}0 + 1{,}0 + 4{,}0 = 168{,}0$. Kolejność dodawania nie ma znaczenia, ważne jest, żeby żadna pozycja nie wypadła.}
  \exitem{Ile godzin z czterdziestu czterech godzin poza snem i pracą zjadają dwie największe kategorie i jaki to procent?}{Gotowanie $8{,}0$ i scrollowanie $7{,}5$ dają $15{,}5$ godziny. Stosunek $15{,}5 / 44 \approx 0{,}35$, czyli około 35\%.}
  \exitem{Zmiana oszczędza 3,0 godziny tygodniowo, kosztuje 0,5 godziny jednorazowo i nie ma kosztu bieżącego. Po ilu dniach się zwraca?}{$n = 0{,}5 / 3{,}0 \approx 0{,}17$ tygodnia, czyli $0{,}17 \cdot 7 \approx 1{,}2$ dnia.}
  \exitem{Zestaw trzech zmian z przykładu ma oszczędność 8,0 godziny tygodniowo, koszt jednorazowy 7,5 godziny i cenę przejścia 2,0 godziny. Ile wynosi zysk netto po dwunastu tygodniach?}{$12 \cdot 8{,}0 - (7{,}5 + 2{,}0) = 96{,}0 - 9{,}5 = 86{,}5$ godziny. Cena przejścia liczy się raz, nie co tydzień, dlatego odejmujesz ją od całości, a nie od każdego tygodnia z osobna.}
  \exitem{Po wycięciu przeglądania w telefonie wolne wraca do czterech godzin, a w arkuszu pojawia się nowy wpis \glqq przeglądanie wiadomości\grqq{} po 3 godziny. Co się stało z odzyskanym czasem i jak to sprawdzić?}{Czas został przesunięty, nie odzyskany. Sprawdzasz to, porównując sumę kategorii ekranów przed i po zmianie oraz sumę wolnego czasu. Jeśli ekrany wróciły do poprzedniego poziomu, zmiana nie zadziałała, niezależnie od tego, jak wyglądała na papierze.}
  \exitem{Dlaczego sen nie powinien znaleźć się na liście kandydatów do cięcia, mimo że zajmuje najwięcej godzin?}{Sen nie jest pozycją w budżecie, który można obniżyć bez kosztu, tylko warunkiem sprawnego funkcjonowania w pozostałych godzinach. Największa liczba godzin nie oznacza największej możliwości oszczędności. W tabeli sen jest punktem odniesienia, a nie pozycją do redukcji.}
\end{cwiczenia}

\zadanie{Przez siedem kolejnych dni prowadź kartkę albo arkusz: zapis co najmniej co godzinę, z kategorią i liczbą minut, bez wpisów zbiorczych poza krótkimi zajęciami poniżej dziesięciu minut. W niedzielę zsumuj wszystko do dziesięciu tysięcy osiemdziesięciu minut tygodnia. Miara sukcesu: suma wpisanych minut różni się od 10\,080 o mniej niż 500 minut (poniżej 5\%), a z tabeli wybierasz jedną zmianę, dla której wpisujesz $S$, $C$ i cenę przejścia, oraz datę, w której ją sprawdzisz po dwóch tygodniach.}
